\documentclass[12pt]{article}

% Page setup
\usepackage{geometry}
\geometry{margin=1in}

% Formatting
\usepackage{setspace}
\usepackage{titlesec}
\titleformat{\section}{\large\bfseries}{\thesection}{1em}{}

% Hyperlinks
\usepackage{hyperref}

% Adjust margins for abstract
\usepackage{changepage}

% Title information
\title{\textbf{Theories of Action, Not Lists of Symptoms:\\Re-centering Long COVID Research on Mechanism, Not Noise}}
\author{Timothy J. VanZeben}
\date{}

\begin{document}

\maketitle

\begin{center}
\textbf{Abstract}
\end{center}

\begin{adjustwidth}{1in}{1in}
\small\itshape
Long COVID presents a diverse array of neurological symptoms including cognitive impairment, fatigue, anhedonia, and dysautonomia. While neuroinflammation has been proposed as a key mechanism, this hypothesis suggests an alternative explanation: selective infection of dopaminergic neurons due to individual genetic variability in surface receptor expression. The result is anatomically specific dopamine disruption, which can explain diverse symptom presentations through damage to distinct neural circuits. This model is constrained, testable, and offers clear clinical implications, providing a unifying theory for a condition that has been obscured by descriptive noise.
\end{adjustwidth}

\section*{Introduction: The Field Has Failed}

Five years into the COVID-19 pandemic, Long COVID research remains mired in symptom cataloging, immune generalities, and vague theories of inflammation and reactivation. The result is a fragmented landscape of weak correlates and unfalsifiable models. Now, as governmental funding is withdrawn and institutional interest wanes, the field faces a stark reality: it failed to deliver a coherent theory of causation.

This failure is not a lack of effort, data, or urgency. It is a failure of epistemic architecture. The research ecosystem pursued descriptive saturation instead of theoretical convergence. It rewarded positive results over constrained hypothesis testing. And it elevated biomarker signals like EBV reactivation or vague immune activation to mechanistic status without disciplined causal framing.

What the field never produced was what science is supposed to generate: a simple, testable theory that explains what is happening and why.

\section*{The Missing Theory: A Dopaminergic Model of Long COVID}

The paper \textit{Dopaminergic Neuron Susceptibility in Long COVID} proposes a direct, anatomical, and genetic model: that SARS-CoV-2 selectively infects dopamine-producing neurons in genetically susceptible individuals. This infection impairs or destroys these neurons, disrupting dopamine signaling across functionally distinct brain pathways. The symptom profile of Long COVID is then a reflection of \textit{which pathways} are disrupted.

\begin{itemize}
    \item \textbf{Nigrostriatal disruption} produces fatigue, exertion intolerance, and motor slowness.
    \item \textbf{Mesolimbic dysfunction} causes anhedonia, motivation loss, and post-viral depression.
    \item \textbf{Mesocortical damage} leads to executive dysfunction, ADHD-like symptoms, and ``brain fog.''
    \item \textbf{Tuberoinfundibular disruption} explains dysautonomia, sleep dysregulation, and metabolic effects.
\end{itemize}

This model does not require post-hoc explanations like EBV reactivation or microclots. It does not rely on inflammation as a primary cause. It maps symptoms to well-characterized neuroanatomy and proposes direct virological and genetic mechanisms. It is constrained. It is falsifiable. It is science.

\section*{Why This Model Succeeds Where Others Fail}

The dopaminergic model breaks from the additive epistemology of the field. It does not accumulate correlates. It narrows the causal hypothesis to a tractable question: does SARS-CoV-2 infect dopamine neurons in genetically predisposed individuals? If so, where? And what are the resulting functional deficits?

It also offers testable predictions:
\begin{itemize}
    \item Neuroimaging should reveal reduced dopamine transporter availability.
    \item Symptom profiles should correlate with affected pathways.
    \item Genetic variants in ACE2, TMPRSS2, DRD2, and DAT1 should be overrepresented in Long COVID cases.
    \item Post-mortem tissue should show dopaminergic neuron damage in key brain regions.
\end{itemize}

This model does not dismiss immune involvement or inflammation. It explains them as secondary responses to primary neurovirological damage.

\section*{The Epistemic Failure of Viral Reactivation Models}

EBV reactivation has emerged as a popular explanation for Long COVID fatigue. But its role is a textbook case of how noise becomes institutionalized. EBV reactivation is common, non-specific, and poorly correlated with symptoms. Its elevation in early studies was a result of weak controls, unstratified samples, and exploratory immunology --- not causal clarity. Its persistence in the literature reflects the field's desperation for mechanisms, not its success in finding them.

The dopaminergic model renders EBV reactivation unnecessary. Fatigue, anhedonia, and dysautonomia all have better explanations via dopamine pathway disruption. The presence or absence of EBV has no bearing on the plausibility of this mechanism. The virus is a confound, not a cause.

\section*{Reclaiming Scientific Rigor}

To move forward, the Long COVID field must abandon its descriptive loop. It must shift from symptom aggregation to theory pruning. The dopaminergic model shows how to do this: propose a tractable mechanism, constrain it anatomically and genetically, and test it directly.

The collapse of funding is not just political. It is epistemic. The field failed to deliver the thing science exists to produce: understanding. This model is an invitation to begin again --- not with more data, but with better theory.

\section*{Citation}

VanZeben, T.J. (2025). \textit{Dopaminergic Neuron Susceptibility in Long COVID: A Hypothesis on Genetic Surface Variability and SARS-CoV-2 Infection.} Independent Research.

\bigskip
\noindent Contact: \texttt{tvanzeben@unm.edu}

\end{document}
